% Include after the abstract in the course-approved ISMIR template.
\section{Methods}
We evaluate five nearest neighbours using cosine distance on standardized track
features. PCA and UMAP maps are evaluated against the original neighbourhoods
using trustworthiness and overlap. We compare k-means and HDBSCAN in feature
space using silhouette, adjusted Rand index against top-level genre, and noise
rate. Silhouette excludes noise; ARI includes noise assignments. Four
one-factor sensitivity checks vary UMAP neighbourhood size and minimum
distance, and HDBSCAN minimum cluster size and minimum samples.

\section{Evaluation results}
% Insert study.json eligible counts and Precision@5 for aligned representations.
% Insert projection_metrics.csv and cluster_metrics.csv results and runtimes.
% Justify selected configurations with scores, coverage and seed sensitivity.
% Describe the audited listening examples with track IDs and audible evidence.
% Report feature extraction time and measured app startup and interaction latency.
% Insert observed task completion, median times and recurring participant feedback.

\section{Limitations}
Genre provides a coarse proxy for perceived similarity. Excluding incompletely
labelled retrieval lists can change which seeds contribute to the reported
mean. A two-dimensional map can distort distances. Runtime depends on hardware,
cache state and compilation. A small convenience sample of classmates limits
generalization of usability observations.
